
\documentclass[]{mcdowellcv}
\usepackage{amsmath}

\name{Michael Raba}
\contacts{(859) 972-4725 \linebreak \href{mailto:michael.raba@uky.edu}{michael.raba@uky.edu} }

\begin{document}
\makeheader


\begin{cvsection}{Featured Projects}
  \begin{cvsubsection}{Projects}{}{}
    \begin{itemize}


\item \textbf{1. Multichamber Muffler Design} (2023). Designed a multichamber muffler system with internal components including chambers, baffles, perforates, and fiberglass absorbents. Simulated acoustic performance using industry-standard tools and validated results against measured data.
  \begin{itemize}
    \item \textbf{Tools Used}: \textbf{ANSYS 2023 (Acoustics Module)}, SIDLAB 5.1, (Verification of Model using 1d approximation)
    \item Modeled transmission and insertion loss across 0–1000 Hz using fluid approximations
    \item Created CAD-style schematics and component breakdowns for internal geometry
    \item Compared simulated and measured acoustic data using Transmission Loss (TL) and Insertion Loss (IL) metrics
    \item Published downloadable simulation files for reproducibility and benchmarking
  \end{itemize}

 \item \textbf{2. Viscoplastic Modeling} (2023). Applied Anand’s viscoplasticity model to simulate strain-rate and temperature-dependent behavior of solder alloys on microchip package. Forward Euler integration scheme in Python to solve constitutive equations and model stress evolution under thermal effects to verify Ansys model.
    \begin{itemize}
      \item \textbf{Tools Used}: \textbf{ANSYS (nonlinear materials module)}, Python NumPy (Verificiation using Numerical Model)
      \item \textbf{Material parameter fitting} based on experimental data
    \end{itemize}


 \item \textbf{3. Electromagnetics Simulation: Finite Difference Time Domain --- Navier Stokes Simulation} (2019). Calculates in 3D space location of a particle in a \textbf{magnetized fluid flow} using the Yee Scheme, using finite difference equations and staggered time-stepping.

\item \textbf{4. MSc Project – POD Analysis of Turbulent Pipe Flow} (2023). Developed a MATLAB-based framework to analyze turbulent rotating pipe flow using Proper Orthogonal Decomposition (POD). Compared Classic and Snapshot POD methods to extract dominant energetic structures and reconstruct flow fields.
  \begin{itemize}
    \item \textbf{Tools Used}: MATLAB, Python, C++, fortran2020, linux
    \item Models produces a PCA simplified model of ~5 mb from 5 Terabytes of data to enable better design decisions
  \end{itemize}
\end{itemize}
  \end{cvsubsection}
\end{cvsection}



\begin{cvsection}{Employment}
  \begin{cvsubsection}{Research Assistant}{University of Kentucky}{2020 -- 2023}
    \begin{itemize}
      \item Collaboration with 11 PhD students at University of Maryland, College Park and University of Oxford
    \end{itemize}
  \end{cvsubsection}

  \begin{cvsubsection}{Teaching Assistant}{University of Kentucky}{Fall 2020 \& Spring 2023}
    \begin{itemize}
      \item TA for course on Design of Experiments (DOE), (ME311 \texttt{Experimental Design II}), a senior level course covering Heat Transfer, Fluid Mechanics, Statics, Engineering Statistics and \texttt{R}; also taught \texttt{ME 330 Fluid Mechanics I} as an undergraduate and \texttt{ME 321 Thermodynamics II} as recitation leader in Fall 2020.
    \end{itemize}
  \end{cvsubsection}

  \begin{cvsubsection}{Wayfinding Assistant}{Chandler Medical Center}{Fall 2019 -- Spring 2020 and July 2023 -- Present}
    \begin{itemize}
      \item Developed custom wayfinding map app to assist hospital visitors. Flask, Python, SQL
    \end{itemize}
  \end{cvsubsection}
\end{cvsection}


\begin{cvsection}{Education}
  \begin{cvsubsection}{Lexington, KY}{University of Kentucky}{Fall 2015 -- December 2023}
    \begin{itemize}
      \item BA in Mathematics, May 2019. In-major GPA: 3.6.
      \item MSc in Mechanical and Aerospace Engineering, December 2023 (Expected). GPA: 3.3
    \end{itemize}
  \end{cvsubsection}
\end{cvsection}

\end{document}
